\emph{(i).} The algebra of this clause is Theorem \claimref{synthesis.canonical-part-of-a-spinor-on-a-null-plane}(iv), where its argument (formerly given here) now stands, with the step that a Lorentz-covariant mass term is a right multiplier computed there: with $\ip{k}{m} = 2$ every vector is $a = \tfrac12 m\,\ip{k}{a} + \tfrac12 k\,\ip{m}{a} + a_{\perp}$, so $k\nabla = \tfrac12 km\,(k\cdot\nabla) + k\nabla_{\perp}$ has no derivative off $N_{k}$; $km = 2 + 2B = 4\Lambda(k)$, $k\Lambda(k) = 0$ and $k$ anticommutes with $\gamma_{1}, \gamma_{2}$, so $k\psi X = k\Lambda(m)\psi X$ for every right multiplier $X$; a mass term commuting with the rotor action $\psi \mapsto R\psi$ is a right multiplication ($\psi \mapsto \psi\gamma_{0}$ for a Dirac mass, Theorem \claimref{field-theory.second-quantisation}; $\psi \mapsto \psi\sigma_{1}$, the charge conjugate, composed with a symmetric matrix on the internal index, for a Majorana mass), together with linear maps on the internal index; so a massive spinor equation $\nabla\psi I_{4}\sigma_{3} = \mathcal{M}(\psi)$, multiplied on the left by $k$, reads $2(k\cdot\nabla)(\Lambda(k)\psi)I_{4}\sigma_{3} = -k\nabla_{\perp}\Lambda(m)\psi\,I_{4}\sigma_{3} + k\,\mathcal{M}(\Lambda(m)\psi)$: along each generator it fixes $\Lambda(k)\psi$ from $\Lambda(m)\psi$ (up to its value at one end of the generator), and the mass enters only here. The flux of the current through $N_{k}$, which defines the canonical part, is $2|\Lambda(m)\psi|^{2}$ by an algebraic identity in $\psi$ that involves no field equation (Theorem \claimref{synthesis.canonical-part-of-a-spinor-on-a-null-plane}(ii); Theorem \claimref{synthesis.light-ray-good-part-of-a-spinor-is-two-chirality-arrays}(i)). So the free data on $N_{k}$ are $\Lambda(m)\psi$, with the same anticommutator, for every mass matrix: one complex canonical channel per Weyl component of the field content (Theorem \claimref{synthesis.light-ray-good-part-of-a-spinor-is-two-chirality-arrays}(iii), (v)). (Two units appear in this step: the spinor values carry Hestenes' unit, right multiplication by $I_{4}\sigma_{3}$, as in the massive equation above, and the numbers of the slot module and of its internal index are read in $\mathrm{span}\{1, I_{4}\}$. They act on different factors and are not identified here; a count of complex channels is half the real dimension of the data, so it is the same for either unit.) (The mass independence of the null-plane data is the familiar property of light-front quantisation, Kogut and Soper 1970, cited from memory as orientation; the argument above does not use it.) The script of this result still checks $k^{2} = 0$, $km = 2 + 2B$, the splitting of each $\gamma_{\mu}$, the anticommutation of $k$ with $\gamma_{1}, \gamma_{2}$, $\mathrm{rank}(L_{k}\Lambda(k)) = 0$, $\mathrm{rank}(L_{k}\Lambda(m)) = 4$, and $\mathrm{rank}(L_{k}R_{X}\Lambda(k)) = 0$ for each of the sixteen basis blades $X$. A Dirac field has two Weyl components and a Majorana field one, so $2n_{D} + n_{W}$ of the light-ray result is the number of Weyl components; how they group into species is decided by the mass matrix and is invisible to the count.

\emph{The link to residue $0$.} $f = \tfrac12(1 + \gamma_{0}\gamma_{3}) = \tfrac12(1 - \gamma_{3}\gamma_{0}) = \Lambda(m)$; $(\gamma_{0}\gamma_{3})^{2} = 1$, so the even algebra of the plane of $k$ and $m$ is $\mathrm{span}\{1, \gamma_{0}\gamma_{3}\} \cong \R \oplus \R$ with the two primitive idempotents $f$ and $1 - f = \Lambda(k)$, $f(1 - f) = 0$, $mf = 0 \neq kf$. This is the single chiral ideal of the residue-$0$ signature $(1,1)$ of Theorem \claimref{synthesis.lowest-causal-residues-zero-four-content}(ii) with $e_{t} = \gamma_{0}$, $e_{a} = \gamma_{3}$; there, mass needs one constituent from each ideal, as here the mass couples $\Lambda(m)\psi$ to $\Lambda(k)\psi$.

\emph{(ii).} Build $V = \Lambda^{\bullet}\C^{5}$ from the empty state by the creation operators of the five modes (the exterior model of the library), with number operators $n_{i} = a^{\dagger}_{i}a_{i}$; on the basis state with occupation $(o_{1}, \dots, o_{5})$ the slot number is $n = o_{1} + o_{2}$, the pair number $k = o_{3} + o_{4} + o_{5}$, $Y = (3n - 2k)/6$, $T_{3} = (o_{1} - o_{2})/2$, and the colour weights are $o_{2 + i} - k/3$, $i = 1, 2, 3$. The gauge algebra preserves $(n, k)$, and the isotypic components are the twelve spaces $F_{n} \otimes \Lambda^{k}W$, of which the two neutral ones, $(n, k) = (0, 0)$ and $(2, 3)$, are both $(1,1)_{0}$, and the other ten form five conjugate pairs: $(0,1)$ with $(2,2)$, $(0,2)$ with $(2,1)$, $(0,3)$ with $(2,0)$, $(1,0)$ with $(1,3)$, $(1,1)$ with $(1,2)$ (the script finds the pairing by comparing weight multisets). By Theorem \claimref{synthesis.spin10-clause-content-is-the-neutral-singlet-choice}(i) a gauge-commuting CPT halving takes one member of each conjugate pair into $V_{+}$ and the other into $V_{-}$, and splits the neutral block into two lines. Each pair contributes the same dimension to both halves, and the neutral block one dimension to each, so $\dim V_{+} = \dim V_{-} = 16$ for all $2^{5}$ sign patterns and every line; the script enumerates the $32$ patterns with each of the two coordinate lines (other lines in the neutral block have the same weights) and finds $\dim V_{\pm} = 16$ and the weights of $V_{-}$ equal to the negatives of those of $V_{+}$ in all $64$ cases. By (i) a generation, the left-handed Weyl fields valued in $V_{+}$, has $16$ complex canonical channels per light ray, whatever the mass matrix; the antiparticles in $V_{-}$ are the holes of the same arrays (the conjugate entry of Theorem \claimref{synthesis.light-ray-good-part-of-a-spinor-is-two-chirality-arrays}(iv)), not further arrays. With one (Q) version 3 array per channel and $\tfrac12\ln 2$ per array (Theorem \claimref{synthesis.half-integer-weight-array-carries-half-ln2}(i)) a generation carries $8\ln 2$, eight registered units of two arrays, and $N_{g}$ generations $8N_{g}$. Removing the neutral singlet leaves fifteen components, an odd number, so not the half of a $32$-dimensional module, and $\tfrac{15}{2}$ units.

\emph{(iii).} In a halving each non-neutral isotypic component of $V$ occurs once, and its conjugate is the other member of its pair, which the halving puts in $V_{-}$; distinct pairs carry disjoint weight multisets. So for a nonzero weight $w$ of $V_{+}$ the weight $-w$ lies only in $V_{-}$, and the zero weight of $V_{+}$ is the neutral line, once. The script confirms this for all $64$ cases, and as a contrast that $V$ itself is vector-like (every weight occurs with its negative), so the test can fail. A Hestenes spinor field valued in a representation $R$ has left-handed part in $R$ and right-handed part in $R$, the latter being a left-handed field in $\bar{R}$; a gauge-covariant Dirac field in the content therefore needs $R \oplus \bar{R}$ inside $V_{+}$, which by the above happens only for $R$ the neutral singlet, and a single neutral singlet pairs only with itself (a Majorana field). Among $N_{g}$ generations the $N_{g}$ neutral singlets give at most $\lfloor N_{g}/2 \rfloor$ such pairs; the script finds no nonzero weight paired with its negative in $N_{g}$ copies, $N_{g} = 1, \dots, 4$. At the anomaly-free halving, $(-1)^{n + k}$ with the particle half containing $(2,3)_{1/6}$ (Theorem \claimref{synthesis.cpt-closure-forces-slot-module-and-chirality}(2)), the hypercharge histogram of $V_{+}$ is $\{0: 1, -\tfrac23: 3, \tfrac16: 6, -\tfrac12: 2, 1: 1, \tfrac13: 3\}$, with sum $0$ and no nonzero value with its negative.

\emph{(iv).} After electroweak breaking the unbroken algebra is colour and the electric charge $T_{3} + Y$. Take the halving $(-1)^{n + k}$ (the clause is stated for $\pm(-1)^{n + k}$ only). Labelling the states of its $V_{+}$ by (charge, colour weights), a Dirac species needs a label and its negative. The labels with nonzero charge pair exactly: charge $\tfrac23$ in colour $c$ with charge $-\tfrac23$ in $\bar{c}$ (three pairs), $-\tfrac13$ with $\tfrac13$ (three), $-1$ with $1$ (one), seven charged Dirac species; the two remaining labels are neutral and colourless (the $\nu$ of the doublet and $\nu^{c}$). With a nonzero $\nu^{c}\nu^{c}$ entry (Theorem \claimref{standard-model.majorana-element}(1)) the neutral mass matrix is a symmetric $2 \times 2$ matrix $\begin{pmatrix} 0 & d \\ d & M \end{pmatrix}$ with $M \neq 0$; its two singular values differ (they coincide only at $M = 0$), so its Takagi eigenvectors are two Majorana fields of different masses, which no unitary mixing combines into one Dirac field, so $n_{D} = 7$, $n_{W} = 2$, $2n_{D} + n_{W} = 16$; with no Majorana entry and a Dirac entry coupling $\nu^{c}$ to the doublet's neutral member, the two neutrals form one Dirac species, $n_{D} = 8$, $n_{W} = 0$, again $16$. In the halving with the other neutral line the existence of such a Dirac entry depends on the bilinears available to the full state and is not decided here; if there is none, the two neutrals are Majorana or massless Weyl fields, $n_{D} = 7$, $n_{W} = 2$, and the count is $16$ again. The count is per four-dimensional field. A five-dimensional field is, level by level, a tower of four-dimensional fields of masses set by its Kaluza-Klein levels, and by (i) each level's Weyl components give one array each on a light ray whatever its mass: one array set per Kaluza-Klein level, equivalently per resolution cell of the fifth direction. The area law fixes only the product of arrays per light ray and cells per Planck area, so this agrees with the common count of the dark-matter premise P1 (one frozen point, one count for every tower) when that count fixes the number of light rays, with tower-dependent arrays per light ray; were every tower to carry the same number of boundary modes instead, the five-dimensional tower's cell density would have to be smaller by its level count. Kaluza-Klein levels counted as emission channels of the exterior are a different count, not constrained here. Without $\nu^{c}$ the charged pairing is unchanged and one neutral component is left, $2 \cdot 7 + 1 = 15$. The script computes these pairings from the weights.

\emph{Conclusion.} The count of arrays per light ray is the number of Weyl components (i), sixteen per generation for every halving (ii); the grouping into Dirac species is neither gauge covariant (iii) nor, under a nonzero Majorana term, complete (iv). Nothing here uses the area law; the residue's calibration, which fixes the product of arrays per light ray and cells per Planck area, is untouched. The result is transport-absent.
